We prove the converse with the stated refinement. Write
$Q : u_{pp}\mathcal{S} \to u_p\mathcal{S}$ for the localization functor.
For $Y = (U, \phi : V \to u(U), y)$ and $b : V' \to V$, put
$$
Y_b = (U, \phi \circ b : V' \to u(U), y),\qquad
k_b = (\text{id}_U, b, \text{id}_y) : Y_b \to Y.
$$
The implication already proved shows that $Q(k_b)$ is strongly cartesian.
Consequently every morphism $h : X \to Y$ over $b$ factors uniquely as
$h = Q(k_b) \circ i_h$, with $i_h : X \to Y_b$ over $\text{id}_{V'}$.
Moreover, $h$ is strongly cartesian if and only if $i_h$ is an isomorphism.
Indeed, if $h$ is strongly cartesian, its universal property gives
$j : Y_b \to X$ over $\text{id}_{V'}$ with $h \circ j = Q(k_b)$.
Uniqueness for $Q(k_b)$ gives $i_h \circ j = \text{id}_{Y_b}$, and
uniqueness for $h$ gives $j \circ i_h = \text{id}_X$.
The reverse implication follows by composition with an isomorphism.

\medskip\noindent
We first establish the required detection of vertical isomorphisms.
Fix $\phi : V \to u(U)$ and a morphism $\delta : x \to y$ in
$\mathcal{S}_U$. Write
$$
D_\delta = (\text{id}_U, \text{id}_V, \delta) :
(U,\phi,x) \longrightarrow (U,\phi,y).
$$
We claim that $Q(D_\delta)$ is invertible if and only if there exist
$c : W \to U$ and $\psi : V \to u(W)$ with $u(c) \circ \psi = \phi$
such that the pullback of $\delta$ along $c$ is invertible.
More explicitly, choose strongly cartesian morphisms
$\gamma_x : x_W \to x$ and $\gamma_y : y_W \to y$ over $c$.
The pullback is the unique vertical morphism $\delta_W : x_W \to y_W$
such that
$$
\gamma_y \circ \delta_W = \delta \circ \gamma_x.
$$
Changing these two cartesian lifts conjugates $\delta_W$ by the unique
vertical comparison isomorphisms, so invertibility is independent of
these choices. If $\delta_W$ is invertible, the displayed equation and
the two morphisms of $R$ defined by $\gamma_x,\gamma_y$ show that
$Q(D_\delta)$ is invertible.

\medskip\noindent
Conversely, represent an inverse to $Q(D_\delta)$ by $Q(e)Q(r)^{-1}$, where
$$
r = (c,\text{id}_V,\gamma) : (T,\theta,z) \to (U,\phi,y),\qquad
e = (a,\text{id}_V,\epsilon) : (T,\theta,z) \to (U,\phi,x),
$$
and $r \in R$. Thus $a,c : T \to U$ and
$u(a)\theta = \phi = u(c)\theta$.
Let $h : T_1 \to T$ be the equalizer of $a,c$.
As $u$ preserves this equalizer, there exists
$\theta_1 : V \to u(T_1)$ with $u(h)\theta_1 = \theta$.
Choose a strongly cartesian $\zeta : z_1 \to z$ over $h$ and refine
the inverse roof by $(h,\text{id}_V,\zeta) \in R$.
Put $d = ah = ch$ and $y_1 = z_1$.
Then $\tau_y = \gamma\zeta : y_1 \to y$ is strongly cartesian over $d$.
Choose a strongly cartesian $\tau_x : x_1 \to x$ over $d$.
There are unique vertical morphisms
$\eta : y_1 \to x_1$ and $\delta_1 : x_1 \to y_1$ satisfying
$$
\tau_x\eta = \epsilon\zeta,\qquad
\tau_y\delta_1 = \delta\tau_x.
$$
For this paragraph, regard $x_1,y_1$ as the triples with base $T_1$
and structure map $\theta_1$, and write
$$
t_x = (d,\text{id}_V,\tau_x),\quad
t_y = (d,\text{id}_V,\tau_y),\quad
E_\eta = (\text{id}_{T_1},\text{id}_V,\eta),\quad
D_1 = (\text{id}_{T_1},\text{id}_V,\delta_1).
$$
The inverse is $Q(t_x)Q(E_\eta)Q(t_y)^{-1}$, whereas
$Q(D_\delta) = Q(t_y)Q(D_1)Q(t_x)^{-1}$.
The two inverse identities therefore give
$$
Q(E_\eta D_1) = \text{id}_{(T_1,\theta_1,x_1)},\qquad
Q(D_1 E_\eta) = \text{id}_{(T_1,\theta_1,y_1)}.
$$
The equality criterion for right fractions, applied to denominators
equal to identities, supplies two arrows of $R$. Write their
$\mathcal{C}$-components as $c_x : T_x \to T_1$ and
$c_y : T_y \to T_1$, their structure maps as $\theta_x,\theta_y$, and
their strongly cartesian $\mathcal{S}$-components as
$\rho_x : z_x \to x_1$ and $\rho_y : z_y \to y_1$.
Then
$$
(\eta\delta_1)\rho_x = \rho_x,\qquad
(\delta_1\eta)\rho_y = \rho_y,\qquad
u(c_x)\theta_x = \theta_1 = u(c_y)\theta_y.
$$
Form $T_2 = T_x \times_{T_1} T_y$, with projections $e_x,e_y$, and put
$k = c_xe_x = c_ye_y$. Preservation of this fibre product gives
$\theta_2 : V \to u(T_2)$ such that
$u(e_x)\theta_2 = \theta_x$ and $u(e_y)\theta_2 = \theta_y$.
In particular $u(k)\theta_2 = \theta_1$ and $u(dk)\theta_2 = \phi$.
Choose strongly cartesian arrows $\kappa_x : x_2 \to x_1$ and
$\kappa_y : y_2 \to y_1$ over $k$, and the unique vertical arrows
$\delta_2 : x_2 \to y_2$ and $\eta_2 : y_2 \to x_2$ with
$$
\kappa_y\delta_2 = \delta_1\kappa_x,\qquad
\kappa_x\eta_2 = \eta\kappa_y.
$$
Since $k = c_xe_x$, cartesianness of $\rho_x$ factors $\kappa_x$
through $\rho_x$, and hence $\eta\delta_1\kappa_x = \kappa_x$.
Similarly $\delta_1\eta\kappa_y = \kappa_y$.
It follows that
$$
\kappa_x\eta_2\delta_2
= \eta\kappa_y\delta_2
= \eta\delta_1\kappa_x
= \kappa_x,
$$
and
$$
\kappa_y\delta_2\eta_2
= \delta_1\kappa_x\eta_2
= \delta_1\eta\kappa_y
= \kappa_y.
$$
Cartesian uniqueness, with fixed identity base maps, yields
$\eta_2\delta_2 = \text{id}_{x_2}$ and
$\delta_2\eta_2 = \text{id}_{y_2}$.
The composites $\tau_x\kappa_x,\tau_y\kappa_y$ are strongly cartesian
over $dk$, and
$$
(\tau_y\kappa_y)\delta_2 = \delta(\tau_x\kappa_x).
$$
Thus $\delta_2$ is the required invertible pullback of $\delta$.
This proves the detection claim without assuming that localization
reflects isomorphisms before a refinement.

\medskip\noindent
Return to $A = (a,b,\alpha) : X_1 \to X_2$.
Choose a strongly cartesian $\tau : a^*x_2 \to x_2$ over $a$, and
write $\alpha = \tau\delta$, where $\delta : x_1 \to a^*x_2$ is vertical.
Put $X_a = (U_1,\phi_1 : V_1 \to u(U_1),a^*x_2)$.
The exact factorization is
$$
A = k_b \circ (a,\text{id}_{V_1},\tau)
\circ (\text{id}_{U_1},\text{id}_{V_1},\delta),
$$
where $k_b : (X_2)_b \to X_2$, and the middle factor belongs to $R$.
Consequently, if $Q(A)$ is strongly cartesian, its vertical factor
relative to $Q(k_b)$ is invertible, and so is $Q(D_\delta)$.
The detection claim supplies $c : W \to U_1$, $\psi : V_1 \to u(W)$
with $u(c)\psi = \phi_1$, strongly cartesian
$\gamma : z \to x_1$ and $\rho : w \to a^*x_2$ over $c$, and an
invertible $\delta_W : z \to w$ such that
$\rho\delta_W = \delta\gamma$.
Therefore
$$
\alpha\gamma = \tau\delta\gamma = \tau\rho\delta_W
$$
is strongly cartesian. The arrow
$r = (c,\text{id}_{V_1},\gamma) : (W,\psi,z) \to X_1$
belongs to $R$, as required.
Conversely, if such an $r$ exists, the implication proved earlier
shows that $Q(Ar)$ is strongly cartesian. Since $Q(r)$ is an
isomorphism, $Q(A) = Q(Ar)Q(r)^{-1}$ is strongly cartesian as well.
This proves the stated characterization.
For any chosen roof $Q(A)Q(r_0)^{-1}$ representing a strongly
cartesian morphism, the same argument gives an additional $r_1 \in R$
with cartesian numerator in the refined roof
$Q(Ar_1)Q(r_0r_1)^{-1}$. Thus the characterization applies to the
roof representation used below, without asserting that an arbitrary
original numerator is already strongly cartesian.
